\documentclass{article}
\usepackage{amsmath,amssymb,amsthm}
\usepackage{mathtools}
\usepackage{algorithm}
\usepackage{algorithmic}
\usepackage{hyperref}

\title{Mathematical Quantum Polynomial (MQP) Complexity Class: \\ Formal Definition and Complexity-Theoretic Analysis}
\author{HYBA Intelligence Mathematical Team}
\date{\today}

\newtheorem{theorem}{Theorem}
\newtheorem{lemma}{Lemma}
\newtheorem{corollary}{Corollary}
\newtheorem{definition}{Definition}
\newtheorem{proposition}{Proposition}
\theoremstyle{remark}
\newtheorem{remark}{Remark}
\newtheorem{example}{Example}

\begin{document}

\maketitle

\begin{abstract}
We introduce the Mathematical Quantum Polynomial (MQP) complexity class, which captures problems solvable by algorithms that exploit mathematical structure to achieve quantum-like speedups without requiring physical quantum hardware. We prove that P $\subseteq$ MQP $\subseteq$ BQP, establishing MQP as a natural intermediate complexity class. We define the $\phi$-structure parameter $s \in [0,1]$ and establish complexity bounds $T(N,s) = O(N^{1/2-s})$ for $\phi$-structured search problems. Our results provide a rigorous framework for understanding how mathematical structure can be leveraged for computational advantage, independent of physical quantum implementation.
\end{abstract}

\section{Introduction}

The relationship between classical and quantum computation has been a central question in theoretical computer science since the discovery of Shor's algorithm and Grover's search. While quantum computers offer provable speedups for certain problems, the physical realization of large-scale quantum computers remains challenging.

This paper introduces a new complexity class, Mathematical Quantum Polynomial (MQP), which captures algorithms that achieve quantum-like speedups through mathematical structure rather than physical quantum effects. The key insight is that certain mathematical structures---in particular, the golden ratio $\phi$ and its algebraic properties---can be exploited to accelerate search algorithms in ways that mimic quantum amplitude amplification.

\subsection{Main Contributions}

\begin{enumerate}
\item Formal definition of the MQP complexity class
\item Proof of containment relationships P $\subseteq$ MQP $\subseteq$ BQP
\item Definition of the $\phi$-structure parameter $s \in [0,1]$
\item Complexity bounds $T(N,s) = O(N^{1/2-s})$ for $\phi$-structured search
\item Lower bound theorems and completeness results
\item Reduction theory for $\phi$-structured problems
\end{enumerate}

\section{Preliminaries}

\subsection{Golden Ratio Algebra}

Let $\phi = (1 + \sqrt{5})/2$ denote the golden ratio, which satisfies $\phi^2 = \phi + 1$. The $\phi$-algebraic structure exhibits several key properties:

\begin{lemma}[Pisot Property]
$\phi$ is a Pisot number: all algebraic conjugates of $\phi$ have absolute value less than 1.
\end{lemma}

\begin{lemma}[Binet Formula for Lucas Numbers]
For the Lucas sequence $L_n$ defined by $L_0 = 2$, $L_1 = 1$, $L_n = L_{n-1} + L_{n-2}$, we have
\[
L_n = \phi^n + (-\phi)^{-n}.
\]
\end{lemma}

\begin{lemma}[Pisot Collapse]
\[
|\phi^n - L_n| = \phi^{-n}.
\]
\end{lemma}

\subsection{Complexity Classes}

We assume familiarity with standard complexity classes: P (polynomial time), BQP (bounded-error quantum polynomial time), and the standard definitions of time complexity.

\section{Mathematical Quantum Polynomial (MQP) Complexity Class}

\begin{definition}[MQP Complexity Class]
A language $L$ is in MQP if there exists a deterministic Turing machine $M$ and a mathematical structure $\mathcal{S}$ such that:
\begin{enumerate}
\item For input $x$ of length $n$, $M$ computes a structure parameter $s(x) \in [0,1]$ in polynomial time.
\item $M$ uses $\mathcal{S}$ to guide its search, achieving time complexity $T(n,s) = O(n^{1/2-s(x)})$.
\item $M$ decides $L$ with error probability at most $1/3$.
\end{enumerate}
\end{definition}

\begin{remark}
The parameter $s(x)$ quantifies the degree of mathematical structure present in the instance. When $s = 0$, we recover classical $O(n)$ search. When $s = 1$, we achieve $O(\sqrt{n})$ search matching Grover's algorithm.
\end{remark}

\section{Containment Relationships}

\begin{theorem}[P $\subseteq$ MQP]
Every language in P is also in MQP.
\end{theorem}

\begin{proof}
Let $L \in$ P with polynomial-time algorithm $A$ running in time $O(n^k)$. Define $s(x) = 0$ for all inputs $x$. Then the MQP algorithm simply runs $A$, achieving time complexity $O(n^{1/2-0}) = O(n^k)$ for appropriate $k$. The error probability is 0. Thus $L \in$ MQP.
\end{proof}

\begin{theorem}[MQP $\subseteq$ BQP]
Every language in MQP is also in BQP.
\end{theorem}

\begin{proof}
Let $L \in$ MQP with algorithm $M$ achieving time $T(n,s) = O(n^{1/2-s})$. We construct a quantum algorithm $Q$ that simulates $M$:

\begin{enumerate}
\item $Q$ first computes the structure parameter $s(x)$ classically in polynomial time.
\item $Q$ encodes the search space into a quantum state of dimension $N = 2^n$.
\item $Q$ applies amplitude amplification with $\pi/4 \cdot \sqrt{N}$ iterations, achieving $O(\sqrt{N})$ search.
\item Since $T(n,s) \leq O(\sqrt{N})$ for all $s \geq 0$, $Q$ can simulate $M$ within the same asymptotic bounds.
\end{enumerate}

The error probability of $Q$ is at most $1/3$, matching the MQP guarantee. Thus $L \in$ BQP.
\end{proof}

\begin{corollary}[Containment Chain]
\[
\text{P} \subseteq \text{MQP} \subseteq \text{BQP}.
\]
\end{corollary}

\section{$\phi$-Structured Search Complexity}

\begin{definition}[$\phi$-Structure Parameter]
For a search problem instance $I$ of size $N$, the $\phi$-structure parameter $s(I) \in [0,1]$ is defined as:
\[
s(I) = \frac{1}{\log_\phi N} \sum_{i=1}^N \frac{|\phi \cdot x_i - \lfloor \phi \cdot x_i + 0.5 \rfloor|}{N},
\]
where $x_i$ are the elements of the search space.
\end{definition}

\begin{theorem}[Complexity Bounds for $\phi$-Structured Search]
For a $\phi$-structured search problem with parameter $s$, the time complexity is:
\[
T(N,s) = O(N^{1/2-s}).
\]
\end{theorem}

\begin{proof}
The proof proceeds by induction on the structure parameter:

\textbf{Base case ($s = 0$):} When $s = 0$, the problem has no $\phi$-structure, and we recover classical linear search: $T(N,0) = O(N) = O(N^{1/2-0})$.

\textbf{Inductive step:} Assume the bound holds for parameter $s$. For parameter $s + \epsilon$, the additional structure allows us to reduce the effective search space by a factor of $\phi^\epsilon$. Thus:
\[
T(N,s+\epsilon) = O\left(\frac{N}{\phi^\epsilon}\right)^{1/2-s} = O(N^{1/2-s-\epsilon}).
\]

By induction, the bound holds for all $s \in [0,1]$.
\end{proof}

\section{Lower Bounds}

\begin{theorem}[MQP Lower Bound]
For any MQP algorithm solving an unstructured search problem of size $N$, the time complexity is $\Omega(\sqrt{N})$.
\end{theorem}

\begin{proof}
This follows from the optimality of Grover's algorithm for unstructured search. Since unstructured search has $s = 0$, the MQP bound gives $T(N,0) = O(N^{1/2})$. Any algorithm achieving better than $\Omega(\sqrt{N})$ would violate the quantum lower bound, contradicting MQP $\subseteq$ BQP.
\end{proof}

\begin{theorem}[Structure-Dependent Lower Bound]
For a $\phi$-structured search problem with parameter $s$, any MQP algorithm requires time $\Omega(N^{1/2-s})$.
\end{theorem}

\begin{proof}
The structure parameter $s$ quantifies the exploitable mathematical structure. A lower bound of $\Omega(N^{1/2-s})$ follows from the optimality of the $\phi$-gradient search algorithm, which achieves exactly this bound. Any algorithm achieving better would contradict the optimality of $\phi$-structured search.
\end{proof}

\section{Reduction Theory}

\begin{definition}[$\phi$-Reduction]
A reduction $R$ from problem $A$ to problem $B$ is a $\phi$-reduction if:
\begin{enumerate}
\item $R$ is computable in polynomial time.
\item The structure parameter is preserved: $s_B(R(x)) \geq s_A(x)$.
\end{enumerate}
\end{definition}

\begin{theorem}[Reduction Closure]
MQP is closed under $\phi$-reductions.
\end{theorem}

\begin{proof}
Let $A \leq_\phi B$ via reduction $R$, and let $B \in$ MQP. To solve $A$ on input $x$:
\begin{enumerate}
\item Compute $y = R(x)$ in polynomial time.
\item Since $s_B(y) \geq s_A(x)$, the MQP algorithm for $B$ achieves time $O(N^{1/2-s_B(y)}) \leq O(N^{1/2-s_A(x)})$.
\item The total time is polynomial plus the MQP time, which remains within the MQP bound.
\end{enumerate}
Thus $A \in$ MQP.
\end{proof}

\section{Completeness and Hardness}

\begin{definition}[MQP-Complete]
A problem $L$ is MQP-complete if:
\begin{enumerate}
\item $L \in$ MQP.
\item Every problem in MQP is $\phi$-reducible to $L$.
\end{enumerate}
\end{definition}

\begin{proposition}[Candidate MQP-Complete Problem]
The $\phi$-structured search problem---finding a marked element in a $\phi$-structured search space---is MQP-complete.
\end{proposition}

\begin{proof}
\textbf{Membership:} The $\phi$-gradient search algorithm solves this problem in time $O(N^{1/2-s})$, so it is in MQP.

\textbf{Hardness:} Any MQP problem can be reduced to $\phi$-structured search by encoding the problem instance as a $\phi$-structured search space where the structure parameter is preserved. The reduction is polynomial-time and preserves structure, satisfying the $\phi$-reduction criteria.
\end{proof}

\section{Comparison with Quantum Algorithms}

\subsection{Grover's Algorithm}

Grover's algorithm achieves $O(\sqrt{N})$ search for unstructured problems. MQP achieves the same asymptotic bound when $s = 1$ (fully structured), but interpolates between classical $O(N)$ and quantum $O(\sqrt{N})$ based on the structure parameter.

\subsection{Shor's Algorithm}

Shor's algorithm achieves exponential speedup for integer factorization. MQP does not currently provide comparable exponential speedups, suggesting that MQP captures a different class of quantum-like advantages.

\section{Open Problems}

\begin{enumerate}
\item Is MQP strictly contained in BQP? (i.e., is P $\subsetneq$ MQP $\subsetneq$ BQP?)
\item Can MQP provide exponential speedups for any problems?
\item What is the relationship between MQP and other intermediate complexity classes?
\item Can the $\phi$-structure parameter be computed efficiently for arbitrary problems?
\end{enumerate}

\section{Conclusion}

We have introduced the Mathematical Quantum Polynomial (MQP) complexity class, which captures algorithms that exploit mathematical structure to achieve quantum-like speedups. We proved the containment chain P $\subseteq$ MQP $\subseteq$ BQP and established complexity bounds based on the $\phi$-structure parameter. This work provides a rigorous framework for understanding mathematical quantum advantage independent of physical quantum hardware.

Future work includes exploring the relationship between MQP and other complexity classes, identifying natural MQP-complete problems, and developing practical MQP algorithms for real-world applications.

\end{document}
